\documentclass[12pt,a4paper]{article}
\usepackage[T2A]{fontenc}
\usepackage[utf8]{inputenc}
\usepackage[russian]{babel}
\usepackage{amsmath,amssymb,amsthm}
\usepackage{geometry}
\usepackage{hyperref}
\usepackage{microtype}
\usepackage{setspace}
\usepackage{caption}
\usepackage{graphicx}

\geometry{left=25mm, right=25mm, top=25mm, bottom=25mm}

\hypersetup{
  colorlinks=true,
  linkcolor=black,
  urlcolor=blue,
  pdfauthor={Зайдиев Артур},
  pdftitle={Отчёт по курсовой работе — Часть №3}
}

\begin{document}
\selectlanguage{russian}

\begin{titlepage}
  \begin{center}
    {Санкт-Петербургский политехнический университет Петра Великого\\
    Институт прикладной математики и механики\\[0.5em]
    \textbf{Высшая школа прикладной математики и вычислительной физики}}

    \vfill

    {\LARGE\bfseries Отчёт по курсовой работе\\ Часть №3\par}
    \vspace{1em}
    {\Large по дисциплине\\[0.3em] {\bfseries «Проект по математической статистике»}}

    \vfill

    \begin{flushright}
        Выполнил\\
        студент гр. 5030102/20202: Зайдиев Артур
    \end{flushright}

    \begin{flushright}
        Проверил\\
        Преподаватель: Заяц Олег Иванович
    \end{flushright}

    \vfill
    Санкт-Петербург\\
    2026
  \end{center}
\end{titlepage}

\setcounter{page}{2}
\newpage

\section*{Вычисление ОИСКО оценки максимального правдоподобия: показательное распределение}

Рассмотрим задачу оценивания плотности показательного распределения методом максимального правдоподобия. Истинная плотность вероятности задана как:
\[
f_0(x) =
\begin{cases}
\lambda e^{-\lambda x}, & x \ge 0, \\
0, & x < 0,
\end{cases}
\]
где параметр распределения $\lambda = 1$.

Пусть $X_1, \dots, X_n$ — независимая выборка из данного распределения. Оценкой максимального правдоподобия (ОМП) параметра $\lambda$ является:
\[
\hat{\lambda}_n = \frac{1}{\bar{X}_n} = \frac{n}{\sum_{i=1}^n X_i},
\]
где $\bar{X}_n = \frac{1}{n}\sum_{i=1}^n X_i$ — выборочное среднее.

Соответствующая оценка плотности вероятности имеет вид:
\[
\hat{f}_n(x) = \hat{\lambda}_n e^{-\hat{\lambda}_n x}, \quad x \ge 0.
\]

\subsection*{Вывод формулы ОИСКО}

Интегральная среднеквадратическая ошибка (ISE) для данной оценки записывается как:
\[
\text{ISE}(\hat{f}_n) = \int_{0}^{\infty} \left( \hat{f}_n(x) - f_0(x) \right)^2 \, dx.
\]

Подставляя выражения для $\hat{f}_n(x)$ и $f_0(x)$, после алгебраических преобразований получаем:
\[
\text{ISE} = (\hat{\lambda}_n - \lambda)^2 \int_{0}^{\infty} x^2 e^{-2\hat{\lambda}_n x} \, dx + \dots
\]

(Полный вывод довольно громоздкий, поэтому обычно переходят к математическому ожиданию.)

Известно, что для показательного распределения случайная величина $\bar{X}_n$ имеет распределение Гамма$(n, n\lambda)$, а $\hat{\lambda}_n$ — обратное гамма-распределение. После вычисления математического ожидания $\mathbb{E}[\text{ISE}]$ и нормировки на $\int f_0^2(x)\,dx = \frac{\lambda}{2}$ получается точная аналитическая формула относительного интегрального среднеквадратического отклонения (ОИСКО):

\[
\bar{\delta}_n = \frac{2n}{(n-1)(n-2)}, \quad n > 2.
\]

Данная оценка обладает **параметрической скоростью сходимости** $O(1/n)$, что асимптотически быстрее, чем у непараметрических методов:
- ядерная оценка $\sim O(n^{-4/5})$,
- гистограммная оценка $\sim O(n^{-2/3})$.

\section*{Построение графиков}

Для сравнительного анализа были построены следующие графики.

\subsection*{График 1. Зависимость нижней границы погрешности $\bar{\delta}_{n,\min}$ (ОМП)}

\begin{figure}[ht!]
    \centering
    \includegraphics[width=0.85\textwidth]{figures/f9.png}
    \caption{ОИСКО оценки максимального правдоподобия $\bar{\delta}_n$ (логарифмический масштаб)}
    \label{fig:omp_delta_min}
\end{figure}

\subsection*{График 2. Сравнение ядерной оценки и ОМП}

\begin{figure}[ht!]
    \centering
    \includegraphics[width=0.85\textwidth]{figures/f10.png}
    \caption{Сравнение ОИСКО ядерной оценки с колокольным ядром и ОМП}
    \label{fig:comparison_kernel_omp}
\end{figure}

\subsection*{График 3. Сравнение гистограммы и ОМП}

\begin{figure}[ht!]
    \centering
    \includegraphics[width=0.85\textwidth]{figures/f11.png}
    \caption{Сравнение ОИСКО гистограммной оценки и ОМП}
    \label{fig:comparison_hist_omp}
\end{figure}

\newpage
\section*{Выводы}

В рамках курсовой работы были исследованы три метода оценивания плотности показательного распределения:
\begin{itemize}
    \item ядерная оценка с колокольным (гауссовским) ядром;
    \item гистограммная оценка;
    \item параметрическая оценка максимального правдоподобия (ОМП).
\end{itemize}

Сравнение методов проводилось по величине относительного интегрального среднеквадратического отклонения $\bar{\delta}_n$.

Оценка максимального правдоподобия обладает параметрической скоростью сходимости $O(1/n)$, в то время как ядерная оценка имеет асимптотическую скорость $O(n^{-4/5})$, а гистограмма — $O(n^{-2/3})$. Поэтому при достаточно больших объёмах выборки ОМП теоретически должна показывать наилучший результат.

Численный анализ подтверждает, что каждый из методов может быть предпочтительным в зависимости от размера выборки. Существуют переходные значения $n$, после которых один метод начинает превосходить другой по точности.

Как правило, наблюдается следующая картина:
\begin{itemize}
    \item При \textbf{малых} объёмах выборки наилучшие результаты демонстрирует \textbf{ядерная оценка} с колокольным ядром.
    \item При \textbf{средних} объёмах выборки часто предпочтительнее оказывается \textbf{гистограммная оценка}.
    \item При \textbf{больших} объёмах выборки преимущество переходит к \textbf{оценке максимального правдоподобия}, что согласуется с её более высокой асимптотической скоростью сходимости.
\end{itemize}

Таким образом, оптимальный выбор метода оценивания плотности существенно зависит от объёма доступных данных. При небольших выборках целесообразно использовать непараметрические методы (в первую очередь ядерную оценку), а при больших — параметрический подход (ОМП).

\end{document}